If $(\xi_j)$ are martingale differences with $|\xi_j|\le d_j$, then for $\theta>0$ convexity of the
exponential between $-d_j$ and $d_j$ and the zero conditional mean give
$\E[\exp(\theta\xi_j)\mid\mathcal F_{j-1}]\le\cosh(\theta d_j)\le\exp(\theta^2d_j^2/2)$. The
compensated product is a positive supermartingale. Stopping it at the first $h$ with
$\sum_{j\le h}\xi_j\ge u$, applying the same calculation to $-\xi_j$, and optimizing over $\theta$ give
\begin{equation}\label{macro:max-mart}
 P\left(\max_{h\le m}\left|\sum_{j\le h}\xi_j\right|\ge u\right)
 \le 2\exp\!\left(-\frac{u^2}{2\sum_{j\le m}d_j^2}\right).
\end{equation}
This is the maximal form of Azuma's inequality~\cite{azuma}. For the fixed partition $\mathcal I$,
there is a $C'_{\mathcal I}$ such that if all the scalar
martingales have absolute maximum less than $b/C'_{\mathcal I}$, then $\mathcal M^*<b$. The factor
$C'_{\mathcal I}$ counts the intervals on a single fiber and the interval and rectangle observations.
The number of tracked lines enters only when we take the union of events. From \eqref{macro:max-mart},
\[
 P^{\rm B}(\mathcal M^*\ge b)\le C_{\mathcal I}(n+1)\exp(-c_{\delta,\mathcal I}b^2 n).
\]

Initially $x=U_f=1$. The lattice sums over intervals and boxes, the approximation of the fiber
endpoints, and the $O(n^3)$ exceptional quadruples give $\Delta(0)\le C_{\mathcal I}/n$, uniformly in
the tracked line. We first choose $\varrho_{\rm mac}$ within the deterministic thresholds, then
$a_{\rm mac}$ and $b$ with $e^{C_\delta}(b+C_\delta a_{\rm mac})<\varrho_{\rm mac}/3$, and finally
take $n$ large with the partition fixed. The discrete Gronwall inequality applied to
\eqref{macro:gronwall} gives $\Delta(s\wedge\tau)<\varrho_{\rm mac}$ whenever $\mathcal M^*<b$,
including the step reaching $\tau$. Thus $\mathcal H_M\subseteq\{\mathcal M^*\ge b\}$, which proves
the bound after absorbing $n+1$ into the exponential.

\emph{Refinement of the partition.} It remains to prove the claim about a fixed observation
partition. Take as $\mathcal I$ a refinement of it of sufficiently small mesh.
Use the common tree of row orders and prefixes of available square choices, stopped at the first
fine exit or at size $M$. At every prefix, each coarse mass and target is the sum of its fine
parts. The freezing time of a line depends only on the chosen squares, so this also holds for
frozen fibers. The triangle inequality, followed for fibers by the maximum over lines, gives
$\Delta_{\rm coarse}\le\Delta_{\rm fine}$ on that same prefix, where the subscripts indicate the
partition with which the discrepancy is computed. A coarse exit at the common radius
is therefore impossible before the first fine exit. If a fine-exit prefix admits a continuation,
any later coarse exit belongs to the event already counted with the full mass of that prefix.
Only transitions from parents strictly before the fine
stop enter this argument, and those parents have a Bethe point by Proposition~\ref{macro:maximo}.
\end{proof}
